\section{SEARCH-R}

\begin{frame}{SEARCH-R: Metadata}
\small
\textbf{SEARCH-R: Structured Entity-Aware Retrieval with Chain-of-Reasoning Navigator for Multi-hop QA}

\paperinfo
{Findings of ACL 2026}
{Dense retrieval plus entity informativeness; ablation mentions \texttt{sentence-transformers/all-*}; NER with \texttt{flair/ner-english-large}}
{Fine-tuned Llama-3.1-8B decomposer; answer generators: Llama-3.1-8B, DeepSeek-V3, GPT-4o-mini; GPT-4o used in analysis}
{Multi-hop question answering on MuSiQue, HotpotQA, and 2WikiMultiHopQA}

\begin{block}{Key insights / contributions}
\scriptsize
\begin{itemize}
  \item Trains an optimized chain-of-reasoning navigator for multi-hop question decomposition.
  \item Combines dense retrieval with entity-aware informativeness scoring from document structure.
  \item Provides more precise multi-document evidence for complex MHQA reasoning.
\end{itemize}
\end{block}
\end{frame}

\begin{frame}{SEARCH-R: Problem and Failure Mode}
\scriptsize
\begin{columns}[T,onlytextwidth]
  \begin{column}{0.48\textwidth}
    \begin{block}{Problem}
      MHQA needs two things at once: a good reasoning path and useful evidence for each hop.
    \end{block}
    \begin{block}{Why prompt-only decomposition fails}
      Large LLMs can generate stochastic or low-utility sub-questions; some paths are logically plausible but do not lead to answerable retrieval.
    \end{block}
  \end{column}
  \begin{column}{0.48\textwidth}
    \begin{block}{Why similarity retrieval fails}
      Dense retrieval may return semantically similar but homogeneous or uninformative documents.
    \end{block}
    \begin{block}{Method response}
      SEARCH-R trains a smaller decomposer using answer-consistency signals and ranks documents using entity informativeness from dependency parse trees.
    \end{block}
  \end{column}
\end{columns}
\end{frame}

\pipelineframe{SEARCH-R}{images/search_r_pipeline.png}{0.98}

\begin{frame}{SEARCH-R: Algorithm Deep Dive}
\scriptsize
\begin{enumerate}
  \item \textbf{Train the navigator:} collect reasoning-path data from strong LLM outputs, filter by answer consistency, then SFT and PPO tune Llama-3.1-8B.
  \item \textbf{Prepare retrieval features:} chunk the corpus, run NER, compute dense embeddings, and compute document entity informativeness from dependency parses.
  \item \textbf{Decompose online:} generate a chain of sub-questions that should expose intermediate entities.
  \item \textbf{Dual retrieval:} retrieve candidates by dense similarity and by entity-aware quantitative information retrieval.
  \item \textbf{Iterative answering:} answer each sub-question, rewrite later sub-questions using previous sub-answers, then synthesize the final response.
\end{enumerate}
\begin{block}{Core intuition}
SEARCH-R improves both the retrieval plan and the evidence scorer: better sub-questions define what to search for, while entity informativeness filters useful evidence.
\end{block}
\end{frame}

\begin{frame}{SEARCH-R: Concrete Example}
\scriptsize
\begin{columns}[T,onlytextwidth]
  \begin{column}{0.48\textwidth}
    \begin{block}{Question}
      What is the experimental satellite that was a forerunner to a communication satellite made by INSAT-4CR's manufacturer?
    \end{block}
    \begin{block}{Reasoning chain}
      \textbf{Step 1:} identify INSAT-4CR's manufacturer.\\
      \textbf{Step 2:} find the manufacturer's communication satellite.\\
      \textbf{Step 3:} retrieve the experimental satellite that preceded it.
    \end{block}
  \end{column}
  \begin{column}{0.48\textwidth}
    \begin{block}{Entity-aware retrieval}
      Dense retrieval asks "which chunk is similar?" SEARCH-R also asks "which entity in this document contributes useful information to the sub-question?"
    \end{block}
  \end{column}
\end{columns}
\end{frame}

\begin{frame}{SEARCH-R: Main Results}
\tiny
\begin{block}{GPT-4o-mini answer generator: F1 comparison}
\centering
\resizebox{0.95\linewidth}{!}{
\begin{tabular}{lccc}
\toprule
Method & MuSiQue & HotpotQA & 2Wiki \\
\midrule
NaiveRAG & 29.82 & 56.92 & 50.61 \\
Iter-RetGen & 38.41 & 57.77 & 58.43 \\
HippoRAG & 46.50 & 56.12 & 62.38 \\
ChainRAG & 50.54 & 59.13 & 56.61 \\
SEARCH-R & \textbf{55.68} & \textbf{60.06} & \textbf{63.86} \\
\bottomrule
\end{tabular}
}
\end{block}
\vspace{0.1cm}
\textbf{Ablation:} on MuSiQue, removing quantitative informativeness retrieval drops F1 from 55.68 to 48.78.
\end{frame}
